% CP-VI-0024
% target_chapter: VI/38 — Projective Morphisms and Closed Embeddings
% source_unit_id: IMP-DL-721195CCB5-PSR004607-99806BF8
% source: Downloads/theory-of-algebraic-geometry-11.tex L4607-L5224
% solution_provenance: SOURCE_IMPLICIT_SOLUTION

\subsection*{Problem 8: Closed Immersion Criterion via Sections}

Let $k$ be an algebraically closed field, let
\[
X\subseteq \mathbb P_k^m
\]
be a closed subscheme, and let
\[
\varphi:X\to \mathbb P_k^n
\]
be a morphism over $\operatorname{Spec}k$ induced by a line bundle
$\mathcal L$ and global sections
\[
s_0,\ldots,s_n\in \Gamma(X,\mathcal L).
\]

Let
\[
V\subseteq \Gamma(X,\mathcal L)
\]
be the $k$-vector subspace spanned by these sections.

Assume that $\mathcal L$ is generated by the sections $s_0,\ldots,s_n$,
so that $\varphi$ is defined everywhere.

We also assume:

\begin{enumerate}
	\item Elements of $V$ separate points.
	That is, for any two distinct closed points $P,Q\in X$, there exists
	$s\in V$ such that
	\[
	s_P\in \mathfrak m_P\mathcal L_P
	\]
	but
	\[
	s_Q\notin \mathfrak m_Q\mathcal L_Q,
	\]
	or vice versa.
	
	\item Elements of $V$ separate tangent vectors.
	That is, for every closed point $P\in X$, the set
	\[
	\{s\in V\mid s_P\in \mathfrak m_P\mathcal L_P\}
	\]
	spans the $k$-vector space
	\[
	\mathfrak m_P\mathcal L_P/\mathfrak m_P^2\mathcal L_P.
	\]
\end{enumerate}

We prove that
\[
\varphi:X\to \mathbb P_k^n
\]
is a closed immersion.

\subsection*{Step 1: Local form of the morphism}

Since the sections $s_0,\ldots,s_n$ generate $\mathcal L$, for every point
$P\in X$ there is at least one index $i$ such that
\[
s_i(P)\neq0.
\]

Let
\[
X_i=\{P\in X\mid s_i(P)\neq0\}.
\]

On $X_i$, the section $s_i$ trivializes $\mathcal L$. Therefore each ratio
\[
\frac{s_j}{s_i}
\]
is a regular function on $X_i$.

On this open set, the morphism $\varphi$ is given by
\[
P\longmapsto
\left[
\frac{s_0}{s_i}(P):
\cdots:
1:
\cdots:
\frac{s_n}{s_i}(P)
\right].
\]

Thus locally the morphism is described by the regular functions
\[
\frac{s_j}{s_i}.
\]

These functions are precisely the pullbacks of affine coordinates on the
standard chart
\[
D_+(y_i)\subseteq \mathbb P_k^n.
\]

\subsection*{Step 2: Point separation implies injectivity}

We show that condition (a) implies that $\varphi$ is injective on closed
points.

Let $P,Q\in X$ be distinct closed points.

By condition (a), there exists $s\in V$ such that, after possibly swapping
$P$ and $Q$,
\[
s_P\in \mathfrak m_P\mathcal L_P
\]
and
\[
s_Q\notin \mathfrak m_Q\mathcal L_Q.
\]

This means that
\[
s(P)=0
\]
but
\[
s(Q)\neq0.
\]

Suppose, for contradiction, that
\[
\varphi(P)=\varphi(Q).
\]

Since $\varphi$ is defined by the projective coordinates
\[
[s_0:\cdots:s_n],
\]
the values of all sections in $V$ at $P$ and $Q$ must have the same
projective vanishing behavior.

In particular, if a section in $V$ vanishes at $P$, then the corresponding
linear form vanishes at $\varphi(P)$, hence also at $\varphi(Q)$, and so the
section vanishes at $Q$.

This contradicts the choice of $s$.

Therefore
\[
\varphi(P)\neq \varphi(Q).
\]

Hence $\varphi$ is injective on closed points.

\subsection*{Step 3: The image is closed}

Since
\[
X\subseteq \mathbb P_k^m
\]
is a closed subscheme of projective space, $X$ is projective over $k$.
In particular, $X$ is proper over $k$.

The target
\[
\mathbb P_k^n
\]
is separated over $k$.

A morphism from a proper scheme to a separated scheme has closed image.
Equivalently, using the fact allowed in the hint, since $X$ is projective,
the image
\[
\varphi(X)
\]
is closed in $\mathbb P_k^n$.

Let
\[
Y=\varphi(X)
\]
with its scheme-theoretic image structure.

Then $\varphi$ factors as
\[
X\xrightarrow{\psi}Y\xhookrightarrow{} \mathbb P_k^n.
\]

To prove that $\varphi$ is a closed immersion, it is enough to show that
\[
\psi:X\to Y
\]
is an isomorphism.

\subsection*{Step 4: Reduction to the structure sheaf}

The map $\psi:X\to Y$ is surjective by definition of $Y$.

By Step 2, it is injective on closed points.

Since $\psi$ is proper and quasi-finite on closed points, and since we are
allowed to use the hint that condition (a) implies that $\varphi$ is a closed
map onto its image, we obtain that $\psi$ is a homeomorphism on underlying
topological spaces.

Thus it remains to show that the morphism of structure sheaves
\[
\mathcal O_Y\longrightarrow \psi_*\mathcal O_X
\]
is an isomorphism.

Since $\psi$ is already a homeomorphism, it is enough to prove that for every
closed point $P\in X$, with
\[
Q=\psi(P)\in Y,
\]
the induced local ring map
\[
\mathcal O_{Y,Q}\longrightarrow \mathcal O_{X,P}
\]
is surjective.

Condition (b) is precisely what gives this surjectivity.

\subsection*{Step 5: Tangent separation and local coordinates}

Fix a closed point
\[
P\in X.
\]

Choose an index $i$ such that
\[
s_i(P)\neq0.
\]

Then $s_i$ trivializes $\mathcal L$ near $P$.

For every section $s\in V$, the quotient
\[
\frac{s}{s_i}
\]
is a regular function near $P$.

Moreover,
\[
s_P\in \mathfrak m_P\mathcal L_P
\]
if and only if
\[
\frac{s}{s_i}\in \mathfrak m_P.
\]

Since $s_i(P)\neq0$, multiplication by $s_i^{-1}$ gives an isomorphism
\[
\mathfrak m_P\mathcal L_P/\mathfrak m_P^2\mathcal L_P
\cong
\mathfrak m_P/\mathfrak m_P^2.
\]

Condition (b) therefore says that the classes of the functions
\[
\frac{s}{s_i},
\qquad
s\in V,\quad s(P)=0,
\]
span the cotangent space
\[
\mathfrak m_P/\mathfrak m_P^2.
\]

But the functions
\[
\frac{s_j}{s_i}
\]
are exactly the pullbacks of affine coordinate functions from the chart
\[
D_+(y_i)\subseteq \mathbb P_k^n.
\]

Hence the image of the local ring
\[
\mathcal O_{Y,Q}
\]
inside
\[
\mathcal O_{X,P}
\]
contains enough functions to generate
\[
\mathfrak m_P/\mathfrak m_P^2.
\]

Equivalently, the induced map on cotangent spaces
\[
\mathfrak m_Q/\mathfrak m_Q^2
\longrightarrow
\mathfrak m_P/\mathfrak m_P^2
\]
is surjective.

Dually, the tangent map
\[
T_PX\longrightarrow T_QY
\]
is injective.

Thus condition (b) says that $\varphi$ does not collapse infinitesimal
directions.

\subsection*{Step 6: Surjectivity of the local ring maps}

Let
\[
A=\mathcal O_{Y,Q},
\qquad
B=\mathcal O_{X,P}.
\]

We have a local homomorphism
\[
A\longrightarrow B.
\]

Its image contains constants and maps the maximal ideal of $A$ into the
maximal ideal of $B$.

From Step 5, the induced map
\[
\mathfrak m_A/\mathfrak m_A^2
\longrightarrow
\mathfrak m_B/\mathfrak m_B^2
\]
is surjective.

Equivalently, the image of $\mathfrak m_A$ generates
\[
\mathfrak m_B
\]
modulo $\mathfrak m_B^2$.

By Nakayama's lemma, the image of $\mathfrak m_A$ generates $\mathfrak m_B$
as a $B$-module.

Since the local functions coming from $Y$ contain these generators of
$\mathfrak m_B$, and since constants also come from $Y$, the local map
\[
\mathcal O_{Y,Q}\longrightarrow \mathcal O_{X,P}
\]
is surjective.

Therefore
\[
\mathcal O_Y\longrightarrow \psi_*\mathcal O_X
\]
is surjective.

Because $\psi$ is already a homeomorphism onto $Y$, this proves that
\[
\psi:X\to Y
\]
is a closed immersion whose underlying topological map is onto $Y$.

Hence $\psi$ is an isomorphism onto $Y$.

\subsection*{Step 7: Conclusion}

We have factored $\varphi$ as
\[
X\xrightarrow{\psi}Y\xhookrightarrow{} \mathbb P_k^n,
\]
where
\[
Y=\varphi(X)
\]
is a closed subscheme of $\mathbb P_k^n$.

We have shown that
\[
\psi:X\to Y
\]
is an isomorphism.

Therefore
\[
\varphi:X\to\mathbb P_k^n
\]
identifies $X$ with the closed subscheme $Y\subseteq\mathbb P_k^n$.

Hence
\[
\boxed{
	\varphi:X\to\mathbb P_k^n
	\text{ is a closed immersion.}
}
\]

\subsection*{Conceptual Summary}

The two hypotheses have clear geometric meanings.

Condition (a) says that the sections separate points:
\[
P\neq Q
\quad\Longrightarrow\quad
\varphi(P)\neq\varphi(Q).
\]

Condition (b) says that the sections separate tangent directions:
\[
\text{no nonzero tangent vector is collapsed by }\varphi.
\]

Together, these imply that $\varphi$ embeds $X$ into projective space with
the correct scheme structure.

Thus
\[
\boxed{
	\text{point separation}+\text{tangent separation}
	\Longrightarrow
	\text{closed immersion}.
}
\]
\section{Concrete Examples: Closed Immersion Criterion via Sections}

\subsection{Example 1: The Identity Embedding of \texorpdfstring{$\mathbb P^1$}{P1}}

Let
\[
X=\mathbb P^1_k
\]
with homogeneous coordinates
\[
[u:v].
\]

Take the line bundle
\[
\mathcal L=\mathcal O_{\mathbb P^1}(1).
\]

The global sections of \(\mathcal O_{\mathbb P^1}(1)\) are spanned by
\[
u,\ v.
\]

Thus
\[
V=\langle u,v\rangle.
\]

These sections define the morphism
\[
\varphi:\mathbb P^1\to\mathbb P^1
\]
given by
\[
[u:v]\longmapsto [u:v].
\]

Thus \(\varphi\) is the identity morphism.

\subsubsection*{Step 1: The sections generate \(\mathcal L\)}

At any point
\[
[u:v]\in\mathbb P^1,
\]
we cannot have both \(u=0\) and \(v=0\).

Therefore at least one of the two sections
\[
u,\ v
\]
is nonzero at every point.

Hence \(u\) and \(v\) generate
\[
\mathcal O_{\mathbb P^1}(1).
\]

Therefore \(\varphi\) is defined everywhere.

\subsubsection*{Step 2: The sections separate points}

Let
\[
P=[a:b],
\qquad
Q=[c:d]
\]
be two distinct points of \(\mathbb P^1\).

Since \(P\neq Q\), the vectors
\[
(a,b)
\qquad\text{and}\qquad
(c,d)
\]
are not scalar multiples.

Consider the section
\[
s=bu-av.
\]

Then
\[
s(P)=ba-ab=0.
\]

Now evaluate at \(Q\):
\[
s(Q)=bc-ad.
\]

If
\[
s(Q)=0,
\]
then
\[
bc-ad=0,
\]
which implies
\[
[a:b]=[c:d].
\]

This contradicts the assumption \(P\neq Q\).

Therefore
\[
s(P)=0,
\qquad
s(Q)\neq0.
\]

So the sections in \(V=\langle u,v\rangle\) separate points.

\subsubsection*{Step 3: The sections separate tangent vectors}

Take the point
\[
P=[1:0].
\]

On the affine chart \(u\neq0\), define the coordinate
\[
t=\frac vu.
\]

Then \(P\) corresponds to
\[
t=0.
\]

The local ring is
\[
\mathcal O_{\mathbb P^1,P}=k[t]_{(t)}.
\]

Its maximal ideal is
\[
\mathfrak m_P=(t).
\]

Thus
\[
\mathfrak m_P/\mathfrak m_P^2
=
(t)/(t^2).
\]

This is a one-dimensional \(k\)-vector space generated by the class of \(t\).

Now we look at sections in \(V\) which vanish at \(P\). Since
\[
u(P)\neq0,
\qquad
v(P)=0,
\]
the section \(v\) vanishes at \(P\).

Trivializing \(\mathcal O_{\mathbb P^1}(1)\) by the nonvanishing section
\(u\), we get
\[
\frac vu=t.
\]

Modulo \((t^2)\), the class of \(t\) generates
\[
\mathfrak m_P/\mathfrak m_P^2.
\]

Therefore the vanishing sections span the cotangent space at \(P\).

The same argument works at every point of \(\mathbb P^1\), after choosing a
chart where one of \(u\) or \(v\) is nonzero.

Hence the sections separate tangent vectors.

\subsubsection*{Conclusion}

The sections \(u,v\) generate \(\mathcal O_{\mathbb P^1}(1)\), separate
points, and separate tangent vectors.

Therefore, by the closed immersion criterion,
\[
\varphi:\mathbb P^1\to\mathbb P^1
\]
is a closed immersion.

Indeed,
\[
\boxed{
	\varphi([u:v])=[u:v]
}
\]
is simply the identity morphism.
